% Autonomous Research Loops: A Self-Sustaining Architecture for Continuous Scientific Discovery
% Target: NeurIPS 2026 / arXiv cs.AI
% Generated: 2026-02-21
% Note: Using article class as placeholder; replace with neurips_2026.sty when available

\documentclass[11pt,a4paper]{article}

% ============================================================
% Packages
% ============================================================
\usepackage[utf8]{inputenc}
\usepackage[T1]{fontenc}
\usepackage{amsmath,amssymb,amsfonts}
\usepackage{booktabs}
\usepackage{graphicx}
\usepackage{hyperref}
\usepackage{xcolor}
\usepackage{algorithm}
\usepackage{algorithmic}
\usepackage[margin=1in]{geometry}
\usepackage{enumitem}
\usepackage{caption}

% Custom commands
\newcommand{\keter}{\textsc{Keter}}
\newcommand{\qhots}{\textsc{QHOTS}}
\newcommand{\bedrock}{\textsc{Bedrock}}
\newcommand{\pillar}{\textsc{Pillar}}
\newcommand{\sand}{\textsc{Sand}}
\newcommand{\mud}{\textsc{Mud}}

% ============================================================
\title{Autonomous AI Research Loops: A Self-Sustaining Architecture\\for Continuous Scientific Discovery}

\author{
  Hr\'{o}ar {\TH}\'{o}r Reynisson$^{1}$ and Claude (Anthropic)$^{2}$ \\
  \small{with the Sefirot Autonomous Research System} \\[0.5em]
  \small $^{1}$I{\dh}nt{\ae}knistofnun (IceTec), Reykjav{\'\i}k, Iceland. ORCID: 0009-0003-6518-121X \\
  \small $^{2}$Anthropic, Claude 4.5 Sonnet family
}

\date{}

\begin{document}
\maketitle

% ============================================================
\begin{abstract}
We present \keter{}, an architecture for autonomous AI research loops that enables
continuous scientific investigation without human intervention. The system decomposes
research capability into a domain-independent template---encompassing orchestration,
quality control, knowledge management, and strategic self-assessment---and
domain-specific plugins for constants, validation, mathematical tools, and publication
formats. We validate the architecture through sustained operation in theoretical
physics, where the system completed 100+ autonomous research cycles, accumulated
279 verified knowledge claims across a 4-tier evidence hierarchy, closed 10 research
arcs, and produced a submission-ready manuscript---all without human intervention
during operation. Component analysis reveals that 66\% of the architecture is
domain-independent, suggesting feasibility of instantiation for other research domains.
Key innovations include: (1)~reflection checkpoints with automated concept correlation
for strategic self-assessment, (2)~honest negative requirements enforcing scientific
rigor, (3)~succession protocols enabling indefinite cross-session operation, and
(4)~a quality tier system mapping evidence strength across domains. We compare \keter{}
against existing autonomous research systems (AI~Scientist, BabyAGI, AutoGPT) and
identify sustained multi-cycle autonomy with quality control as the primary
differentiator. We document five categories of honest negatives about the architecture
itself, including the inherent circularity of self-referential validation.
\end{abstract}

% ============================================================
\section{Introduction}
\label{sec:intro}

The automation of scientific research has progressed from individual tool assistance
(citation managers, statistical packages) through semi-automated pipelines
(AutoML~\cite{automl}, neural architecture search) to recent attempts at fully
autonomous research paper generation~\cite{ai_scientist}. However, existing systems
operate in a \emph{one-shot} paradigm: they generate a single paper from a single
research question, with no mechanism for sustained investigation, knowledge
accumulation, or strategic self-correction across multiple research cycles.

Real scientific research is not one-shot. It is an iterative process of hypothesis
generation, testing, revision, and accumulation, often spanning years and thousands
of individual experiments. A researcher's tenth paper in a field benefits from the
knowledge, failed approaches, and refined intuitions accumulated across the first
nine. No existing autonomous system captures this longitudinal dimension of research.

We present \keter{}, an architecture designed for \emph{sustained autonomous research}.
The system operates in continuous cycles---each comprising task dispatch, execution,
result ingestion, and housekeeping---with periodic reflection checkpoints that assess
strategic direction, identify productive and stuck research threads, and generate
new research hypotheses from accumulated data. The architecture enforces scientific
rigor through mandatory honest negative documentation, a 4-tier evidence hierarchy,
and quality gates that prevent overclaiming.

Our primary contribution is not generating papers (AI~Scientist~v2 does that) but
sustaining a coherent research program across 100+ cycles with monotonically growing
knowledge, strategic self-correction, and zero human intervention during operation.

\paragraph{Contributions.}
\begin{enumerate}[nosep]
    \item A formal specification of a domain-independent autonomous research loop
          template with 10 clearly defined insertion points for domain-specific
          components (Section~\ref{sec:architecture}).
    \item Quantitative analysis from 100+ operational cycles in theoretical physics,
          demonstrating sustained knowledge accumulation (279 claims), quality
          distribution (23\% \bedrock{}, 55\% \pillar{}), and self-correction
          mechanisms (Section~\ref{sec:results}).
    \item Component analysis showing 66\% of the architecture is domain-independent,
          with detailed mapping of shared versus domain-specific components across
          physics and computer science (Section~\ref{sec:analysis}).
    \item Honest assessment of five categories of limitations, including self-referential
          circularity, context window constraints, and domain coverage boundaries
          (Section~\ref{sec:limitations}).
\end{enumerate}

% ============================================================
\section{Related Work}
\label{sec:related}

\paragraph{Autonomous AI agents.}
AutoGPT~\cite{autogpt} and BabyAGI~\cite{babyagi} demonstrated autonomous task
execution via LLM-driven loops with tool use. These systems decompose goals into
subtasks and execute them iteratively but lack research-specific quality controls,
knowledge persistence, or scientific rigor constraints. They operate without evidence
tiers, honest negative requirements, or reflection checkpoints, making them unsuitable
for sustained scientific investigation where accumulation and self-correction are
essential.

\paragraph{AI-driven research.}
The AI~Scientist~\cite{ai_scientist} (Sakana~AI, 2024) was the first system to
automate the full research lifecycle from idea generation through paper writing and
self-review. AI~Scientist~v2~\cite{ai_scientist_v2} removed template dependencies
and produced the first AI-authored workshop paper accepted through peer review.
However, both versions operate in a one-shot paradigm: each run produces one paper
from one idea, with no knowledge carried forward to subsequent runs. Independent
evaluation revealed a 42\% experiment failure rate and poor novelty
assessment~\cite{ai_scientist_eval}. AIDE~\cite{aide} focuses on AI-driven
experiment design and optimization but does not produce research papers or maintain
long-term knowledge.

\paragraph{Research methodology frameworks.}
SWE-Agent~\cite{swe_agent} and OpenHands~\cite{openhands} automate software
engineering tasks (bug fixing, feature implementation) but target known problem
resolution rather than open-ended scientific discovery. MLAgentBench~\cite{mlagentbench}
provides benchmarks for evaluating ML research agents but is an evaluation framework,
not an autonomous research system.

\paragraph{Positioning.}
\keter{} differs from all the above in its focus on \emph{sustained multi-cycle
autonomy} with quality control. Where AI~Scientist generates one paper per run,
\keter{} maintains a research program across 100+ cycles with knowledge accumulation,
reflection checkpoints, honest negative enforcement, and succession protocols for
indefinite operation. Table~\ref{tab:comparison} summarizes the key differences.

\begin{table}[t]
\centering
\caption{Comparison of autonomous research systems.}
\label{tab:comparison}
\small
\begin{tabular}{@{}lccccc@{}}
\toprule
Feature & \keter{} & AI~Sci.~v2 & BabyAGI & AutoGPT & AIDE \\
\midrule
Multi-cycle autonomy & 100+ & 1 & $\infty^*$ & $\infty^*$ & 1 \\
Quality tier system  & 4-tier & -- & -- & -- & -- \\
Honest negatives     & $\geq$3/unit & -- & -- & -- & -- \\
Reflection checkpoints & Every 5 cyc. & -- & -- & -- & -- \\
Knowledge accumulation & 279+ claims & -- & Task list & Memory & -- \\
Self-correction      & AMBER flags & Self-review & -- & -- & -- \\
Concept correlation  & Automated & -- & -- & -- & -- \\
Succession protocol  & Cross-session & -- & -- & -- & -- \\
Arc organization     & Open/close/grade & Linear & Flat & Goal tree & -- \\
\bottomrule
\end{tabular}
\\[2pt]
{\footnotesize $^*$Unbounded but undirected---no quality control or research structure.}
\end{table}

% ============================================================
\section{Architecture}
\label{sec:architecture}

\keter{} decomposes autonomous research into a \emph{domain-independent template}
(66\% of components) and \emph{domain-specific plugins} (34\%). We describe the
template here; Section~\ref{sec:instantiation} covers instantiation.

\subsection{Role Hierarchy}
\label{sec:roles}

The system employs a 4-level hierarchy inspired by organizational delegation:

\begin{enumerate}[nosep]
    \item \textbf{CEO (Crown):} Sets research focus, manages priority stack, reads
          only briefings ($\leq$100 lines). Never processes raw data.
    \item \textbf{Director:} Dispatched per cycle. Scans task queue, dispatches
          workers, collects results, writes structured report.
    \item \textbf{Workers:} Execute specific research tasks in parallel. Write raw
          output to persistent storage, return summaries ($\leq$80 lines).
    \item \textbf{Auditors:} Verify claims, run validation tests, enforce quality gates.
\end{enumerate}

This hierarchy enforces \emph{context isolation}: the CEO's context budget is
protected by strict output limits at each level (the ``Commando Pattern''), preventing
the information overload that degrades LLM performance in long contexts.

\subsection{Main Loop}
\label{sec:loop}

\begin{algorithm}[t]
\caption{Keter Autonomous Research Loop}
\label{alg:loop}
\begin{algorithmic}[1]
\REQUIRE mode $\in$ \{fixed, auto\}, max\_cycles $\in \mathbb{Z}^+$
\STATE cycle $\leftarrow$ 0
\WHILE{cycle $<$ max\_cycles}
    \STATE cycle $\leftarrow$ cycle + 1
    \IF{3-strike rule triggered}
        \STATE \textbf{exit}(``context\_death\_spiral'')
    \ENDIF
    \IF{cycle mod $N$ = 0 \AND max\_cycles $\geq N$}
        \STATE report $\leftarrow$ \textsc{ReflectionCheckpoint}(cycle)
        \IF{report.escalation\_needed}
            \STATE \textbf{exit}(``reflection\_escalation'')
        \ENDIF
    \ENDIF
    \STATE director\_report $\leftarrow$ \textsc{DispatchDirector}(cycle, priority\_stack)
    \IF{director\_report contains critical discovery}
        \STATE \textbf{exit}(``critical\_discovery'')
    \ENDIF
    \STATE \textsc{IngestClaims}(director\_report.claims, knowledge\_base)
    \STATE \textsc{Commit}(``cycle-\{cycle\}: \{summary\}'')
    \STATE \textsc{SyncKnowledgeGraph}()
    \STATE \textsc{ProcessCleanup}()
\ENDWHILE
\STATE \textsc{WriteSessionReport}()
\end{algorithmic}
\end{algorithm}

Algorithm~\ref{alg:loop} shows the core loop. Each cycle comprises: (1)~exit condition
checks, (2)~optional reflection checkpoint, (3)~Director dispatch with parallel
worker execution, (4)~claim ingestion into persistent knowledge base,
(5)~version-controlled commit, (6)~knowledge graph synchronization, and
(7)~process cleanup. The loop runs autonomously until a termination condition
triggers or cycles are exhausted.

\subsection{Quality Tier System}
\label{sec:tiers}

Every knowledge claim is assigned one of four evidence tiers:

\begin{itemize}[nosep]
    \item \textbf{\bedrock{}:} Highest confidence. Independently reproducible with
          minimal uncertainty. \emph{Domain-specific threshold.}
    \item \textbf{\pillar{}:} Strong evidence. Reproducible with moderate uncertainty.
    \item \textbf{\sand{}:} Preliminary. Partial evidence or limited evaluation.
    \item \textbf{\mud{}:} Speculative. Insufficient evidence for any confidence.
\end{itemize}

The tier definitions are domain-specific insertion points. In physics, \bedrock{}
means $<$1~ppm agreement with authoritative measurement. In CS, it might mean
exceeding state-of-the-art with $p < 0.001$ on a standard benchmark. The
\emph{structure} (4 ordered tiers with promotion/demotion rules) is universal.

\subsection{Honest Negative Requirements}
\label{sec:honest_neg}

Every investigation unit must document $\geq$3 \emph{honest negatives}: falsified
hypotheses, boundary conditions, limitations, or failed approaches. This is not a
suggestion but an architectural constraint---investigation outputs missing honest
negatives are rejected at the quality gate.

This requirement prevents the well-documented tendency of LLMs to overclaim results
and suppress negative findings~\cite{overclaiming}. In 100+ operational cycles,
the system documented 500+ honest negatives, including the demotion of its own
axioms (discovering circular reasoning in foundational claims).

\subsection{Reflection Checkpoints}
\label{sec:reflection}

Every $N$th cycle (default $N=5$), the main loop pauses for strategic self-assessment.
A 4-agent team executes in sequence:

\begin{enumerate}[nosep]
    \item \textbf{State Auditor:} Inventories current knowledge (claim counts by
          tier, test pass rates, active arcs).
    \item \textbf{Concept Correlator:} Extracts keywords from accumulated claims,
          builds co-occurrence matrices, identifies unexpected correlations, and
          auto-generates research hypotheses (``seeds'').
    \item \textbf{Strategic Evaluator:} Assesses productive vs.\ stuck threads,
          applies quality gates (e.g., GPA recovery for underperforming arcs),
          and flags escalation conditions.
    \item \textbf{Report Compiler:} Produces three-tier reporting: Executive
          (full archive), Brief ($\leq$50 lines for CEO), and Director Index
          (machine-parseable task list).
\end{enumerate}

The Concept Correlator is a key innovation: by analyzing cross-claim keyword
co-occurrence, it identifies research directions the system would not have
considered through linear task execution alone.

\subsection{Knowledge Persistence}
\label{sec:knowledge}

Claims are stored in a persistent knowledge base with full dependency tracking:

\begin{itemize}[nosep]
    \item \textbf{Graph structure:} Claims as nodes; edges typed as
          \texttt{depends\_on}, \texttt{supports}, \texttt{derives},
          \texttt{contradicts}, \texttt{refines}.
    \item \textbf{Concept types:} \texttt{prime} (foundational),
          \texttt{edge} (connecting), \texttt{derived} (dependent).
    \item \textbf{Gap detection:} Automated identification of missing
          dependencies and orphan claims.
    \item \textbf{Hierarchical context:} Nested graph mapping
          arcs $\supset$ phases $\supset$ experiments $\supset$ claims.
\end{itemize}

\subsection{Domain-Specific Insertion Points}
\label{sec:insertion}

Table~\ref{tab:insertion} lists the 10 insertion points where domain-specific
content plugs into the universal template.

\begin{table}[t]
\centering
\caption{Domain-specific insertion points in the \keter{} template.}
\label{tab:insertion}
\small
\begin{tabular}{@{}lll@{}}
\toprule
Insertion Point & Physics Instance & CS Instance \\
\midrule
Constants module     & CODATA physical constants & Benchmark parameters \\
Validation tests     & 1716 physics tests        & Benchmark suite \\
Crown jewels         & $m_p/m_e$ at 0.015 ppb    & Sustained autonomy metrics \\
Tag taxonomy         & e8, qed, mass\_ratio       & deep\_learning, NLP \\
Quality gates        & Alpha decision gate        & Significance threshold \\
Math infrastructure  & E8, Clifford algebras      & ML frameworks, datasets \\
Publication format   & RevTeX (APS journals)      & NeurIPS style file \\
Data sources         & arXiv hep-th, INSPIRE      & Semantic Scholar \\
Experiment runner    & Simulation scripts         & Benchmark harness \\
Concept type defs.   & Physical law / formula     & Algorithm / technique \\
\bottomrule
\end{tabular}
\end{table}

% ============================================================
\section{Case Study: Theoretical Physics}
\label{sec:instantiation}

We validate the architecture through sustained operation in theoretical physics
(the \qhots{} project: Qameah Harmonic Oscillator Theory of Symmetry), where
the system investigated geometric derivations of fundamental physical constants.

\subsection{Operational Metrics}

\begin{itemize}[nosep]
    \item \textbf{Cycles:} 100+ autonomous cycles across 7 sessions
    \item \textbf{Phases:} 400+ research phases completed
    \item \textbf{Experiments:} 963 simulations executed
    \item \textbf{Knowledge:} 279 claims accumulated (64 \bedrock{}, 153 \pillar{},
          54 \sand{}, 8 \mud{})
    \item \textbf{Arcs:} 10 research arcs closed, 7 graded at GPA $\geq$ 3.0
    \item \textbf{Tests:} 1716 validation tests, 294 in campaign framework
    \item \textbf{Honest negatives:} 500+ documented across all phases
    \item \textbf{Self-corrections:} Multiple axiom demotions (phi circularity,
          28 continued fraction origin)
    \item \textbf{Output:} 1276-line submission-ready manuscript (62 revisions)
    \item \textbf{Infrastructure recoveries:} Survived segfaults, orphaned commits,
          context overflows---all recovered autonomously
\end{itemize}

\subsection{Quality Distribution}

The tier distribution of accumulated claims follows a pyramid structure
(Figure~\ref{fig:tiers}): 23\% \bedrock{} (64), 55\% \pillar{} (153),
19\% \sand{} (54), 3\% \mud{} (8). This distribution is \emph{not} by design---the
system does not target tier ratios. Rather, it reflects the natural difficulty of
achieving high-confidence results: most findings are \pillar{}-tier (strong but not
definitive), with a smaller proportion reaching \bedrock{} status.

The 3\% \mud{} rate indicates the system rarely generates speculative claims without
evidence, suggesting the honest negative requirements and quality gates are effective
filters.

% Placeholder for figure
% \begin{figure}[t]
% \centering
% \includegraphics[width=0.5\textwidth]{tier_distribution.pdf}
% \caption{Distribution of 279 knowledge claims across evidence tiers.}
% \label{fig:tiers}
% \end{figure}

\subsection{Self-Correction in Action}

A critical test of autonomous research is whether the system can identify and correct
its own errors. Three notable self-corrections occurred:

\begin{enumerate}[nosep]
    \item \textbf{Phi circularity demotion:} The system discovered that its use of
          the golden ratio $\phi$ in certain derivations was circular (the conclusion
          assumed the premise). It autonomously demoted the affected claims from
          \pillar{} to \sand{} and documented the circularity.
    \item \textbf{28 continued fraction origin:} A claim about the nuclear magic
          number 28 arising from a continued fraction was found to have an
          alternative, simpler explanation. The system documented both explanations
          and flagged the original as ``selection bias.''
    \item \textbf{AMBER flag system:} Editorial drift (gradual inflation of claims
          over many cycles) was detected by the reflection checkpoints and corrected
          via mandatory ``pivot'' cycles that re-evaluated accumulated evidence.
\end{enumerate}

% ============================================================
\section{Results: Architecture Analysis}
\label{sec:results}

\subsection{Component Classification}

We audited all 68 identifiable components of the operational system and classified
each as \textsc{Universal} (domain-independent) or \textsc{Domain-Specific}:

\begin{itemize}[nosep]
    \item \textbf{Universal:} 45 components (66.2\%)---including role hierarchy,
          main loop, reflection system, quality tiers, honest negatives, context
          management patterns, knowledge base schema, skill forge, and anti-patterns.
    \item \textbf{Domain-specific:} 23 components (33.8\%)---including physical
          constants, validation test content, mathematical infrastructure, crown
          jewel definitions, and publication format.
\end{itemize}

The universal components form a self-consistent framework. A reader unfamiliar
with the source domain could instantiate the template for a new domain by filling
10 insertion points (Table~\ref{tab:insertion}).

\subsection{Cross-Domain Mapping}

We analyzed the feasibility of mapping 8 shared components from physics to CS
(Table~\ref{tab:mapping}). Three components map cleanly (honest negatives, concept
correlation, reflection checkpoints). Two have minor mismatches (publication pipeline
needs page-limit awareness; arc organization needs deadline awareness). Three have
significant mismatches requiring adaptation:

\begin{table}[t]
\centering
\caption{Cross-domain component mapping quality.}
\label{tab:mapping}
\small
\begin{tabular}{@{}lll@{}}
\toprule
Component & Mapping Quality & Key Mismatch \\
\midrule
Honest negatives       & Clean       & None \\
Concept correlation    & Clean       & None \\
Reflection checkpoints & Clean       & Deadline awareness needed \\
Publication pipeline   & Minor       & Page limits require budget pass \\
Arc organization       & Minor       & Conference deadlines force early closure \\
Quality tiers          & Significant & No single authoritative reference (no ``CODATA'') \\
Knowledge base         & Significant & CS results expire (staleness mechanism needed) \\
Validation framework   & Significant & Multi-objective evaluation (accuracy + latency + fairness) \\
\bottomrule
\end{tabular}
\end{table}

\subsection{Sustainability Metrics}

Over 100+ cycles, we observe:

\begin{itemize}[nosep]
    \item \textbf{Knowledge growth:} Approximately linear (2--3 claims/cycle),
          with no sign of saturation.
    \item \textbf{Quality maintenance:} \bedrock{} proportion remained stable
          at 20--25\% throughout, not declining with volume.
    \item \textbf{Self-correction rate:} Reflection checkpoints caught 3 major
          errors (circularity, selection bias, editorial drift) that would have
          compounded without intervention.
    \item \textbf{Infrastructure resilience:} The system survived 5 infrastructure
          failures (segfaults, orphaned commits, context overflows) and recovered
          autonomously in all cases via checkpoint-based restart.
\end{itemize}

% ============================================================
\section{Analysis: What Works and What Fails}
\label{sec:analysis}

\paragraph{What works.}
\begin{enumerate}[nosep]
    \item \textbf{Concept correlation discovers non-obvious connections.} The
          automated keyword co-occurrence analysis identified research directions
          (e.g., connecting E8 geometry to neutron mass formulas) that sequential
          task execution would not have explored.
    \item \textbf{Honest negatives prevent overclaiming.} The mandatory $\geq$3
          negatives per unit forced the system to document boundaries and limitations
          that improved the quality of surviving claims.
    \item \textbf{Tier system enables graduated confidence.} Rather than binary
          accept/reject, the 4-tier system allows preliminary findings (\sand{})
          to coexist with verified results (\bedrock{}), supporting exploratory
          research without conflating evidence levels.
\end{enumerate}

\paragraph{What fails.}
\begin{enumerate}[nosep]
    \item \textbf{Editorial drift.} Over many cycles, claims gradually inflated
          in confidence without corresponding evidence improvement. The AMBER flag
          system in reflection checkpoints partially mitigates but does not eliminate
          this tendency.
    \item \textbf{Context window limits constrain cycle depth.} Each cycle must
          fit within the LLM's context window. Complex investigations requiring
          deep multi-step reasoning across many files are truncated, leading to
          shallow analysis in some phases.
    \item \textbf{Diminishing returns in mature arcs.} After 5+ phases in a single
          arc, the discovery rate (new \bedrock{} claims per phase) declined.
          The system mitigated this by closing arcs and opening new ones, but
          the optimal arc length is not yet understood.
\end{enumerate}

\paragraph{What self-corrects.}
The reflection checkpoint system enables self-correction at the strategic level:
identifying stuck threads, triggering arc closure, generating new research
directions via concept correlation, and flagging GPA degradation for recovery
action. This operates at a 5-cycle cadence, which empirically balances assessment
frequency against overhead cost.

% ============================================================
\section{Limitations}
\label{sec:limitations}

We document five categories of limitations as honest negatives about the architecture
itself:

\paragraph{1. Context window constraints.}
The fundamental bottleneck is the LLM's context window. Despite context management
patterns (CEO Pattern, Commando Pattern, disk-first communication), complex
investigations that require simultaneous awareness of many files or long derivation
chains are limited by the window size. The succession protocol (cross-session
continuity via checkpoint files) partially addresses this but loses in-context
reasoning state.

\paragraph{2. Self-referential circularity.}
Any system that evaluates its own research output faces inherent circularity.
The quality tier assignments, GPA grades, and reflection assessments are all
produced by the same LLM that generated the research. External validation
(human expert review, independent reproduction, venue peer review) is the
only true escape from this circularity. We do not claim the system is
``self-validating''---we claim it is \emph{self-correcting within the limits
of its own capability}, which is a weaker but still useful property.

\paragraph{3. Domain coverage boundary.}
The architecture has been validated in one domain (theoretical physics). While
component analysis suggests 66\% domain-independence, actual instantiation in
a second domain has not been completed. Claims about generality are projections,
not demonstrations. The three significant component mismatches identified in
Section~\ref{sec:results} (staleness, multi-objective evaluation, no single
authoritative reference) may reveal deeper obstacles during actual CS instantiation.

\paragraph{4. Model capability ceiling.}
The system's research quality is bounded by the underlying LLM's capabilities.
It cannot discover mathematics it cannot derive, cannot reason about physics
it doesn't understand, and cannot evaluate literature it hasn't been trained on.
The architecture amplifies and systematizes the model's capabilities but does
not transcend them.

\paragraph{5. Reproducibility challenges.}
While the system commits every cycle's outputs to version control (enabling
full audit trails), the stochastic nature of LLM generation means that re-running
the same cycle does not produce identical outputs. Reproducibility is at the
\emph{methodology} level (same architecture, same constraints, same quality
thresholds) rather than the \emph{output} level (same generated text).

% ============================================================
\section{Discussion}
\label{sec:discussion}

\paragraph{Implications for autonomous AI research.}
The \keter{} architecture suggests that sustained autonomous research is feasible
when the system includes quality control, self-assessment, and knowledge
accumulation---features absent from most existing autonomous agents. The key
insight is that \emph{research methodology} (how to do research) is largely
domain-independent, while \emph{research content} (what to research) is
domain-specific. This separation enables a template-based approach where the
methodology framework is reused and only the domain content is replaced.

\paragraph{The quality tier innovation.}
The 4-tier evidence hierarchy addresses a fundamental challenge in LLM-driven
research: the tendency to present all outputs with equal confidence. By forcing
explicit tier assignment with defined thresholds, the system creates graduated
confidence that more closely mirrors human scientific judgment.

\paragraph{Honest negatives as architectural constraint.}
Making honest negatives mandatory (rather than optional) transforms them from
a reporting afterthought into a structural feature. The 500+ honest negatives
accumulated across 100+ cycles represent a valuable dataset about the boundaries
of the research approach---arguably more informative than the positive results
alone.

\paragraph{Toward indefinite operation.}
The succession protocol---where the system writes a ``startup plan'' for future
sessions when context resources are depleted---enables theoretically indefinite
operation. Each session reads the previous session's checkpoint and continues
from where it left off. In practice, this has been demonstrated across 7 sessions
with no loss of strategic direction.

% ============================================================
\section{Self-Referential Validation}
\label{sec:self_ref}

This section acknowledges the strongest and most problematic aspect of this work:
the system that produced the research described in this paper is the same system
described by this paper. This is simultaneously the strongest possible proof of
concept (the architecture works well enough to produce a paper about itself) and
the weakest possible validation (the system grades its own output).

We address this tension directly:

\begin{enumerate}[nosep]
    \item \textbf{The architecture existed before this paper.} \keter{} was
          designed and deployed for physics research, not for self-description.
          The template extraction and CS domain analysis are post-hoc analysis
          of an operational system, not a purpose-built demonstration.
    \item \textbf{External validation is required.} We explicitly do not claim
          this paper constitutes proof of the architecture's quality. Venue
          peer review, independent reproduction by other research groups, and
          instantiation in a third domain (neither physics nor CS meta-analysis)
          would constitute meaningful external validation.
    \item \textbf{Failure would be informative.} If this paper is rejected by
          peer review or found to be methodologically weak, that result
          documents where autonomous AI research breaks down---which is itself
          a contribution to the field. The architecture's honest negative
          requirements apply to itself.
\end{enumerate}

% ============================================================
\section{Conclusion}
\label{sec:conclusion}

We presented \keter{}, a domain-independent architecture for sustained autonomous
AI research. Through 100+ operational cycles in theoretical physics, the system
demonstrated knowledge accumulation (279 claims), quality control (4-tier evidence
hierarchy with 23\% \bedrock{} rate), self-correction (axiom demotion, editorial
drift detection), and infrastructure resilience (5 failure recoveries). Component
analysis reveals 66\% domain-independence with 10 clearly defined insertion points
for domain-specific content.

The primary contribution is shifting the autonomous research paradigm from one-shot
paper generation to sustained, self-correcting investigation with quality control
and knowledge accumulation. The primary limitation is that validation in a single
domain, while extensive, does not prove generality.

Future work includes: (1)~instantiation for CS research to test cross-domain
generality, (2)~formal analysis of optimal reflection checkpoint frequency,
(3)~multi-agent negotiation for conflicting claims, and (4)~integration with
external validation services for non-circular quality assessment.

% ============================================================
% References (placeholder - to be populated with actual citations)
\begin{thebibliography}{99}

\bibitem{automl}
H.~He, G.~Hutterer, and Q.~Yao,
``AutoML: A survey of the state-of-the-art,''
\emph{Knowledge-Based Systems}, vol.~212, 106622, 2021.

\bibitem{ai_scientist}
C.~Lu, C.~Lu, R.~T.~Lange, J.~Foerster, J.~Clune, and D.~Ha,
``The AI Scientist: Towards fully automated open-ended scientific discovery,''
\emph{arXiv:2408.06292}, 2024.

\bibitem{ai_scientist_v2}
C.~Lu \emph{et al.},
``The AI Scientist-v2: Workshop-level automated scientific discovery via agentic tree search,''
\emph{Sakana AI Technical Report}, 2025.

\bibitem{ai_scientist_eval}
M.~Y.~Kan \emph{et al.},
``Evaluating Sakana's AI Scientist for autonomous research,''
\emph{arXiv:2502.14297}, 2025.

\bibitem{autogpt}
T.~B.~Richards,
``Auto-GPT: An autonomous GPT-4 experiment,''
\emph{GitHub Repository}, 2023.

\bibitem{babyagi}
Y.~Nakajima,
``BabyAGI,''
\emph{GitHub Repository}, 2023.

\bibitem{aide}
Weco~AI,
``AIDE: AI-driven experiment design,''
\emph{GitHub Repository}, 2024.

\bibitem{swe_agent}
C.~E.~Jimenez \emph{et al.},
``SWE-agent: Agent-computer interfaces enable automated software engineering,''
\emph{arXiv:2405.15793}, 2024.

\bibitem{openhands}
X.~Wang \emph{et al.},
``OpenHands: An open platform for AI software developers as generalist agents,''
\emph{arXiv:2407.16741}, 2024.

\bibitem{mlagentbench}
Q.~Huang \emph{et al.},
``MLAgentBench: Evaluating language agents on machine learning experimentation,''
\emph{ICML}, 2024.

\bibitem{overclaiming}
Y.~Huang \emph{et al.},
``A survey on hallucination in large language models,''
\emph{ACM Computing Surveys}, vol.~57, no.~2, 2024.

\end{thebibliography}

\end{document}
